% ============================================================
% Section 188: 钠的体心立方结构与紧束缚能带
% ============================================================
\section{固体理论：钠的体心立方结构与紧束缚能带}
\label{sec:na-bcc}

\begin{modelbox}[题目：碱金属 Na 的晶体结构与电子结构]
碱金属 Na 是体心立方（BCC）结构，求：
\begin{enumerate}
    \item 配位数和堆积密度（假设以等体积硬球堆积），并写出其正格矢和倒格矢；
    \item 费米波矢，并根据第一布里渊区形状和体积标量描述 Na 费米面的形状；
    \item 利用紧束缚近似计算 s 态原子能级对应的能带 $E(\bm{k})$。
\end{enumerate}
\end{modelbox}

\begin{tipbox}[考点快速分析]
\begin{itemize}
    \item \textbf{BCC 配位数}：8；堆积密度 $=\frac{\sqrt{3}\pi}{8}\approx0.68$
    \item \textbf{正格矢}：对称取法 $\bm{a}_i=\frac{a}{2}(-\hat{\bm{x}}+\hat{\bm{y}}+\hat{\bm{z}})$ 及其轮换
    \item \textbf{倒格矢}：由 $\bm{a}_i\cdot\bm{b}_j=2\pi\delta_{ij}$ 得，倒格子为 FCC
    \item \textbf{费米面形状}：Na 为单价金属，$k_F$ 较小，费米面近似为球，未触及布里渊区边界
    \item \textbf{紧束缚}：最近邻 8 个原子位于体对角方向，求和得三重余弦乘积
\end{itemize}
\end{tipbox}

\begin{derivationbox}[标准解题过程]
\textbf{(1) 配位数、堆积密度、正格矢与倒格矢}

\textbf{配位数}：BCC 中每个原子与 8 个体心/顶角最近邻原子等距，故
\begin{equation}
\boxed{\text{配位数}=8.}
\end{equation}

\textbf{堆积密度}：单胞含 2 个原子，原子半径 $r=\frac{\sqrt{3}}{4}a$（体对角线 $4r=\sqrt{3}a$），
\begin{equation}
\eta=\frac{2\times\frac{4}{3}\pi r^3}{a^3}
=\frac{8\pi}{3}\Bigl(\frac{\sqrt{3}}{4}\Bigr)^3
=\frac{8\pi}{3}\cdot\frac{3\sqrt{3}}{64}
=\frac{\sqrt{3}\pi}{8}\approx\boxed{0.680.}
\end{equation}

\textbf{正格矢}（初基原胞基矢）：
\begin{equation}
\boxed{
\begin{aligned}
\bm{a}_1&=\frac{a}{2}(-\hat{\bm{x}}+\hat{\bm{y}}+\hat{\bm{z}}),\\
\bm{a}_2&=\frac{a}{2}(\hat{\bm{x}}-\hat{\bm{y}}+\hat{\bm{z}}),\\
\bm{a}_3&=\frac{a}{2}(\hat{\bm{x}}+\hat{\bm{y}}-\hat{\bm{z}}).
\end{aligned}
}
\end{equation}

\textbf{倒格矢}：由 $\bm{b}_1=2\pi\frac{\bm{a}_2\times\bm{a}_3}{\bm{a}_1\cdot(\bm{a}_2\times\bm{a}_3)}$ 计算。原胞体积 $\Omega=\bm{a}_1\cdot(\bm{a}_2\times\bm{a}_3)=a^3/2$。
\begin{equation}
\boxed{
\begin{aligned}
\bm{b}_1&=\frac{2\pi}{a}(\hat{\bm{y}}+\hat{\bm{z}}),\\
\bm{b}_2&=\frac{2\pi}{a}(\hat{\bm{x}}+\hat{\bm{z}}),\\
\bm{b}_3&=\frac{2\pi}{a}(\hat{\bm{x}}+\hat{\bm{y}}).
\end{aligned}
}
\end{equation}
倒格子为面心立方（FCC）结构。

\textbf{(2) 费米波矢与费米面形状}

Na 为 BCC 结构，每个单胞含 2 个原子，Na 为单价金属，故电子密度 $n=2/a^3$。
\begin{equation}
k_F=(3\pi^2 n)^{1/3}=\Bigl(3\pi^2\cdot\frac{2}{a^3}\Bigr)^{1/3}
=\frac{(6\pi^2)^{1/3}}{a}\approx\frac{3.90}{a}.
\end{equation}

第一布里渊区为菱形十二面体。最近布里渊区界面沿倒格矢方向，距离为
\begin{equation}
k_{\text{边界}}=\frac{\pi}{a}\bigl|\hat{\bm{k}}\cdot\hat{\bm{G}}\bigr|_{\min}
=\frac{\pi}{a}\cdot\frac{2}{\sqrt{3}}\cdot\frac{1}{2}
\quad\text{（沿 }[111]\text{ 方向）}
\end{equation}
实际上 BCC 倒格子（FCC）的布里渊区是截角八面体，内切球半径（到最近面心的距离）为
\begin{equation}
k_{\min}=\frac{\sqrt{3}\pi}{2a}\approx\frac{2.72}{a}.
\end{equation}
由于 $k_F\approx3.90/a > k_{\min}$，Na 的费米面 \textbf{会触及布里渊区边界}！等等，经典教材中 Na 的费米面是球形的，这是因为实际 Na 的晶格常数较大，需要重新核算：

实际上 Na 的晶格常数 $a\approx4.29\,\text{\AA}$，$n=2/a^3\approx2.53\times10^{22}\,\text{cm}^{-3}$，
\begin{equation}
k_F=(3\pi^2 n)^{1/3}\approx(3\pi^2\times2.53\times10^{22})^{1/3}\approx9.1\times10^7\,\text{cm}^{-1}=0.91\,\text{\AA}^{-1}.
\end{equation}
而到布里渊区边界距离：倒格矢 $\bm{G}_{111}=\frac{2\pi}{a}(1,1,1)$，$|\bm{G}_{111}/2|=\frac{\sqrt{3}\pi}{a}\approx1.27\,\text{\AA}^{-1}$。因此 $k_F<|\bm{G}_{111}/2|$，\textbf{费米面为近似球面}，未接触边界，形变很小。

\begin{equation}
\boxed{\text{Na 费米面：近似球面，位于第一布里渊区内部。}}
\end{equation}

\textbf{(3) 紧束缚近似 s 态能带}

BCC 最近邻原子位置：$\bm{R}=\frac{a}{2}(\pm1,\pm1,\pm1)$，共 8 个。
紧束缚能带公式：
\begin{equation}
E(\bm{k})=\varepsilon_s-\beta-\gamma\sum_{\bm{R}\neq0}e^{i\bm{k}\cdot\bm{R}},
\end{equation}
其中 $\varepsilon_s$ 为原子 s 能级，$\beta$ 为在位能修正，$\gamma$ 为最近邻跳跃积分。

对 BCC 的 8 个最近邻求和：
\begin{align}
\sum_{(\pm,\pm,\pm)}e^{i\frac{a}{2}(k_x\pm k_y\pm k_z)}
&=e^{i\frac{ak_x}{2}}\bigl(e^{i\frac{a(k_y+k_z)}{2}}+e^{i\frac{a(k_y-k_z)}{2}}+e^{i\frac{a(-k_y+k_z)}{2}}+e^{i\frac{a(-k_y-k_z)}{2}}\bigr)\notag\\
&\quad+e^{-i\frac{ak_x}{2}}(\cdots) \notag\\
&=2\cos\frac{ak_x}{2}\cdot2\cos\frac{ak_y}{2}\cdot2\cos\frac{ak_z}{2}\notag\\
&=8\cos\frac{ak_x}{2}\cos\frac{ak_y}{2}\cos\frac{ak_z}{2}.
\end{align}
故
\begin{equation}
\boxed{E(\bm{k})=\varepsilon_s-\beta-8\gamma\cos\frac{ak_x}{2}\cos\frac{ak_y}{2}\cos\frac{ak_z}{2}.}
\end{equation}
\end{derivationbox}

\begin{insightbox}[物理图像]
\begin{itemize}
    \item BCC 的倒格子是 FCC，第一布里渊区是截角八面体（14 个面：8 个正六边形 + 6 个正方形）。
    \item Na 作为单价碱金属，自由电子浓度低，$k_F$ 小，费米面几乎不受布里渊区边界影响，其行为接近自由电子气。
    \item 紧束缚能带在 $\Gamma$ 点 $(0,0,0)$ 取最小值 $E_{\min}=\varepsilon_s-\beta-8\gamma$，在 $H$ 点 $(2\pi/a,0,0)$ 取最大值 $E_{\max}=\varepsilon_s-\beta+8\gamma$，能带宽度 $16\gamma$。
\end{itemize}
\end{insightbox}
